% v11 RQ1/RQ2 on independent human-annotated, multi-language data (WP C, protocol frozen before evaluation).
\section{RQ1: Independent Benchmark and Student Code}\label{sec:external}
\subsection{Human-Annotated Multi-Language Benchmark}
Our syntax-tree reference shares its construction with syntax-aware differencers and is not human ground truth. We therefore evaluated every method on the public CodeMapper benchmark \cite{hu2026codemapper}: \cExtHumN{} code regions whose later locations were annotated by people in two sets (one with manually and one with automatically chosen source regions), plus tool-derived Java histories and a suppression set, \cExtExtN{} regions in all. We ran CodeMapper's pinned artifact unchanged and GumTree's official release as additional comparators. The protocol, adapters and frozen guard code were fixed before any of these data were evaluated, so the results are held out for every method. Two limits apply. Living Review sees only the pair of files in each case, not a whole project, so cross-file recovery is not exercised. And our fresh CodeMapper runs did not reproduce its published numbers exactly (for example \cRepAObserved{} against \cRepAPublished{} exact regions of \cRepAN{} in the first set); we report our own runs.

% Generated by tools/v11/build_wpc_v11_tex.py; do not edit.
\begin{table}[t]
\caption{Human-Annotated (Independent) and Multi-Language Benchmarks}\label{tab:external}
\centering\footnotesize\setlength{\tabcolsep}{2.5pt}
\begin{tabular}{@{}lrrrrrr@{}}\toprule
 & \multicolumn{4}{c}{Human A/B (200)} & Non-Py & All \\ \cmidrule(lr){2-5}
Method & Line & Chars & Missed & Silent & Line & Line \\ \midrule
Deployed cascade & 66.0 & 24.5 & 8.5 & 3.5 & 63.7 & 57.5 \\
All guards, $\theta=.85$ & 70.0 & 26.0 & 1.5 & 0.5 & 67.8 & 64.5 \\
All guards, $\theta=.55$ & 77.0 & 17.5 & 1.5 & 0.5 & 77.4 & 70.9 \\
Cascade v1 & 70.0 & 21.5 & 8.5 & 3.5 & 67.1 & 59.5 \\
Contextual tracker & 85.5 & 28.0 & 3.8 & 0.0 & 86.3 & 70.2 \\
GumTree & 62.0 & 17.0 & 1.5 & 0.5 & 71.2 & 66.8 \\
CodeMapper & 91.5 & 74.0 & 4.6 & 1.0 & 89.7 & 86.2 \\
\bottomrule
\end{tabular}
\tablenote{}{Percentages. Line: file and start line, including correct absence. Chars: exact character region. All guards: IG, the original copy guard AG and VW (Table~\ref{tab:locator} uses the corrected AGi). CodeMapper: this evaluation's run and change adapter (Table~\ref{tab:locator} uses a later run with the cascade's quotation labels). Non-Py: the 146 human-annotated regions outside Python. All: all 2,392 external regions, including tool-derived Java histories. Living Review sees only the paired files. GumTree execution failures count as incorrect. The pinned CodeMapper is not fully deterministic (Section~\ref{sec:prefixes}).}
\end{table}


On the human-annotated regions the deployed cascade found the correct line, or correctly reported the region gone, for \cExtHumDepAcc\%, against \cExtHumTrkAcc\% for the contextual tracker and \cExtHumCmAcc\% for CodeMapper (Table~\ref{tab:external}). The gap is larger for exact regions: CodeMapper recovered \cExtHumCmChars\% of the annotated character regions exactly, the cascade \cExtHumDepChars\%. Outside Python the cascade reached \cExtHumNpDepAcc\% (CodeMapper \cExtHumNpCmAcc\%). On the full external set the gap persisted (\cExtExtDepAcc\% against \cExtExtCmAcc\%). CodeMapper, a region mapper built for this task, is substantially more accurate at locating annotated code than the cascade, and the cascade's fixed-height windows recover region boundaries poorly.

Safety is a different question. The deployed cascade left \cExtHumDepMiss\% of the needed re-checks unrequested on the human-annotated regions and placed \cExtHumDepSilent\% of all regions on wrong code without a request. With the three guards in the version evaluated here (IG, the original copy guard AG and VW; Section~\ref{sec:locator} uses the corrected AGi) these fell to \cExtHumAllMiss\% and \cExtHumAllSilent\%, and accuracy rose to \cExtHumAllAcc\% (\cExtHumAllCalAcc\% at the calibrated threshold). The price is paid on preserved code: re-check requests for the \cExtHumDepFalseN{} regions whose content was unchanged rose from \cExtHumDepFalse\% to \cExtHumAllFalse\%. CodeMapper missed \cExtHumCmMiss\% of the re-checks under this evaluation's change adapter, which compares its mapped text with the original; it has no review-state model of its own. Section~\ref{sec:locator} reports a later run of the same pinned CodeMapper with a read-only path resolver and the cascade's own quotation labels, which missed \wiHumCmQMissK{} of \wiHumCmQMissN{} (\wiHumCmQMiss\%). These are small samples (\cExtHumN{} regions), and the human annotations were made for region tracking, not for review concerns. They support two conclusions: the cascade is not a competitive locator, and its value lies in the state contract it feeds; a deployment could replace the matcher with a region mapper such as CodeMapper and keep the same labels, guards and issue lifecycle.
